\chapter{2012年重庆卷（文）}

\section{单选题}

\begin{problem}
命题“若 \(p\) 则 \(q\)”的逆命题是（\quad）

\choices
    {若 \(q\) 则 \(p\)}
    {若 \(\lnot p\) 则 \(\lnot q\)}
    {若 \(\lnot q\) 则 \(\lnot p\)}
    {若 \(p\) 则 \(\lnot q\)}

\begin{answer}
A
\end{answer}

\begin{solution}
逆命题是把原命题的条件与结论互换. 原命题的条件为 \(p\)，结论为 \(q\)，所以其逆命题为“若 \(q\) 则 \(p\)”，故选 A.
\end{solution}
\end{problem}

\begin{problem}
不等式 \(\frac{x-1}{x+2}<0\) 的解集为（\quad）

\choices
    {\((1,+\infty)\)}
    {\((-\infty,-2)\)}
    {\((-2,1)\)}
    {\((-\infty,-2)\cup(1,+\infty)\)}

\begin{answer}
C
\end{answer}

\begin{solution}
分式有意义的条件是 \(x\ne-2\)，并且分子、分母的零点分别为 \(1\) 和 \(-2\). 在三个区间上判断符号：当 \(x\in(-\infty,-2)\) 时，分子、分母均为负，分式为正；当 \(x\in(-2,1)\) 时，分子为负而分母为正，分式为负；当 \(x\in(1,+\infty)\) 时，分式为正. 因此不等式的解集为 \((-2,1)\)，故选 C.
\end{solution}
\end{problem}

\begin{problem}
设 \(A\)，\(B\) 为直线 \(y=x\) 与圆 \(x^2+y^2=1\) 的两个交点，则 \(|AB|=\)（\quad）

\choices
    {\(1\)}
    {\(\sqrt{2}\)}
    {\(\sqrt{3}\)}
    {\(2\)}

\begin{answer}
D
\end{answer}

\begin{solution}
圆 \(x^2+y^2=1\) 的圆心为原点 \(O\)，半径为 \(1\). 直线 \(y=x\) 经过圆心，所以它截圆所得的弦 \(AB\) 是直径，故 \(|AB|=2\times1=2\)，选 D.
\end{solution}
\end{problem}

\begin{problem}
\((1-3x)^5\) 的展开式中 \(x^3\) 的系数为（\quad）

\choices
    {\(-270\)}
    {\(-90\)}
    {\(90\)}
    {\(270\)}

\begin{answer}
A
\end{answer}

\begin{solution}
二项式展开的通项为 \(\mathrm C_5^k(1)^{5-k}(-3x)^k\). 含 \(x^3\) 的项对应 \(k=3\)，其系数为 \(\mathrm C_5^3(-3)^3=10\times(-27)=-270\)，故选 A.
\end{solution}
\end{problem}

\begin{problem}
\(\frac{\sin 47^{\circ}-\sin 17^{\circ}\cos 30^{\circ}}{\cos 17^{\circ}}=\)（\quad）

\choices
    {\(-\frac{\sqrt{3}}{2}\)}
    {\(-\frac{1}{2}\)}
    {\(\frac{1}{2}\)}
    {\(\frac{\sqrt{3}}{2}\)}

\begin{answer}
C
\end{answer}

\begin{solution}
由和角公式，\(\sin47^{\circ}=\sin(17^{\circ}+30^{\circ})=\sin17^{\circ}\cos30^{\circ}+\cos17^{\circ}\sin30^{\circ}\). 因此原式等于 \(\frac{\cos17^{\circ}\sin30^{\circ}}{\cos17^{\circ}}=\sin30^{\circ}=\frac{1}{2}\)，故选 C.
\end{solution}
\end{problem}

\begin{problem}
设 \(x\in\R\)，向量 \(\bs{a}=(x,1)\)，\(\bs{b}=(1,-2)\)，且 \(\bs{a}\perp\bs{b}\)，则 \(|\bs{a}+\bs{b}|=\)（\quad）

\choices
    {\(\sqrt{5}\)}
    {\(\sqrt{10}\)}
    {\(2\sqrt{5}\)}
    {\(10\)}

\begin{answer}
B
\end{answer}

\begin{solution}
由 \(\bs{a}\perp\bs{b}\)，得 \(\bs{a}\cdot\bs{b}=0\)，即 \(x\times1+1\times(-2)=0\)，所以 \(x=2\). 于是 \(\bs{a}+\bs{b}=(3,-1)\)，从而 \(|\bs{a}+\bs{b}|=\sqrt{3^2+(-1)^2}=\sqrt{10}\)，故选 B.
\end{solution}
\end{problem}

\begin{problem}
已知 \(a=\log_{2}3+\log_{2}\sqrt{3}\)，\(b=\log_{2}9-\log_{2}\sqrt{3}\)，\(c=\log_{3}2\)，则 \(a\)，\(b\)，\(c\) 的大小关系是（\quad）

\choices
    {\(a=b<c\)}
    {\(a=b>c\)}
    {\(a<b<c\)}
    {\(a>b>c\)}

\begin{answer}
B
\end{answer}

\begin{solution}
设 \(t=\log_2 3\). 因为 \(3>2\)，所以 \(t>1\). 利用对数的运算性质，\(a=\log_2(3\sqrt3)=\log_2(3^{3/2})=\frac32t\)，\(b=\log_2\frac{9}{\sqrt3}=\log_2(3^{3/2})=\frac32t\)，故 \(a=b\). 又 \(c=\log_3 2=\frac{1}{\log_2 3}=\frac1t\)，且 \(t>1\) 时 \(\frac32t>\frac1t\)，所以 \(a=b>c\)，故选 B.
\end{solution}
\end{problem}

\begin{problem}
设函数 \(f(x)\) 在 \(\R\) 上可导，其导函数为 \(f'(x)\)，且函数 \(f(x)\) 在 \(x=-2\) 处取得极小值，则函数 \(y=xf'(x)\) 的图象可能是（\quad）

\choices
    {\choicebitmap[width=0.15\paperwidth]{img/2012/chongqing_liberal/q8_optA.png}}
    {\choicebitmap[width=0.15\paperwidth]{img/2012/chongqing_liberal/q8_optB.png}}
    {\choicebitmap[width=0.15\paperwidth]{img/2012/chongqing_liberal/q8_optC.png}}
    {\choicebitmap[width=0.15\paperwidth]{img/2012/chongqing_liberal/q8_optD.png}}

\begin{answer}
C
\end{answer}

\begin{solution}
要说明“可能是”，只需构造一个满足条件的可导函数. 取 \(f(x)=\frac12(x+2)^2\)，则 \(f(x)\) 在 \(x=-2\) 处取得极小值，且 \(f'(x)=x+2\). 于是 \(y=xf'(x)=x(x+2)=x^2+2x\)，其图象是开口向上的抛物线，过 \(x=-2\) 和 \(x=0\) 两个零点，在两零点之间位于 \(x\) 轴下方，正是选项 C 所表示的图象，故选 C.
\end{solution}
\end{problem}

\begin{problem}
设四面体的六条棱的长分别为 \(1\)，\(1\)，\(1\)，\(1\)，\(\sqrt{2}\) 和 \(a\)，且长为 \(a\) 的棱与长为 \(\sqrt{2}\) 的棱异面，则 \(a\) 的取值范围是（\quad）

\choices
    {\((0,\sqrt{2})\)}
    {\((0,\sqrt{3})\)}
    {\((1,\sqrt{2})\)}
    {\((1,\sqrt{3})\)}

\begin{answer}
A
\end{answer}

\begin{solution}
设长为 \(a\) 的一条棱为 \(AB\)，与之异面的长为 \(\sqrt2\) 的棱为 \(CD\)，则另外四条棱 \(AC\)，\(AD\)，\(BC\)，\(BD\) 的长均为 \(1\). 取 \(AB\) 的中点为原点，使 \(A=(-\frac a2,0,0)\)，\(B=(\frac a2,0,0)\). 由 \(AC=BC=1\) 及 \(AD=BD=1\)，可设 \(C=(0,u,v)\)，\(D=(0,s,t)\)，并有 \(u^2+v^2=s^2+t^2=1-\frac{a^2}{4}\). 记 \(r=\sqrt{1-\frac{a^2}{4}}\)，则 \(C\)，\(D\) 位于同一圆心为原点、半径为 \(r\) 的圆上.

因为 \(CD=\sqrt2\)，由三角不等式得 \(\sqrt2=|CD|\leq|C|+|D|=2r\)，所以 \(2\leq4-a^2\)，即 \(a^2\leq2\). 若 \(a=\sqrt2\)，则 \(CD=2r\)，等号成立时 \(D=-C\)，四点 \(A\)，\(B\)，\(C\)，\(D\) 共面，不能构成非退化四面体，故必须有 \(a<\sqrt2\). 又棱长 \(a>0\)，所以 \(0<a<\sqrt2\).

反过来，任取 \(0<a<\sqrt2\)，则 \(2r>\sqrt2\)，可在该圆上取两点 \(C\)，\(D\)，使 \(CD=\sqrt2\) 且 \(D\ne-C\)，从而四点不共面并满足六条棱的长度条件. 因此 \(a\) 的取值范围为 \((0,\sqrt2)\)，故选 A.
\end{solution}
\end{problem}

\begin{problem}
设函数 \(f(x)=x^2-4x+3\)，\(g(x)=3x-2\)，集合 \(M=\{x\in\R\mid f(g(x))>0\}\)，\(N=\{x\in\R\mid g(x)<2\}\)，则 \(M\cap N\) 为（\quad）

\choices
    {\((1,+\infty)\)}
    {\((0,1)\)}
    {\((-1,1)\)}
    {\((-\infty,1)\)}

\begin{answer}
D
\end{answer}

\begin{solution}
先求 \(f(g(x))\)：\[
f(g(x))=(3x-2)^2-4(3x-2)+3=9x^2-24x+15=3(3x-5)(x-1).
\]
所以 \(f(g(x))>0\) 的解集为 \(x<1\) 或 \(x>\frac53\). 另一方面，\(g(x)<2\) 等价于 \(3x-2<2\)，即 \(x<\frac43\). 取交集时，区间 \((\frac53,+\infty)\) 与 \((-\infty,\frac43)\) 无交集，因此 \(M\cap N=(-\infty,1)\)，故选 D.
\end{solution}
\end{problem}

\section{填空题}

\begin{problem}
首项为 1，公比为 2 的等比数列的前 \(4\) 项和 \(S_{4}=\fillinblank{}\).

\begin{answer}
\(15\)
\end{answer}

\begin{solution}
首项为 \(1\)，公比为 \(2\)，所以前四项依次为 \(1\)，\(2\)，\(4\)，\(8\). 因此 \(S_4=1+2+4+8=15\).
\end{solution}
\end{problem}

\begin{problem}
若 \(f(x)=(x+a)(x-4)\) 为偶函数，则实数 \(a=\fillinblank{}\).

\begin{answer}
\(4\)
\end{answer}

\begin{solution}
展开得 \(f(x)=x^2+(a-4)x-4a\). 偶函数满足 \(f(-x)=f(x)\)，二次式中一次项系数必须为零，所以 \(a-4=0\)，即 \(a=4\).
\end{solution}
\end{problem}

\begin{problem}
设 \(\triangle ABC\) 的内角 \(A\)，\(B\)，\(C\) 的对边分别为 \(a\)，\(b\)，\(c\)，且 \(a=1\)，\(b=2\)，\(\cos C=\frac{1}{4}\)，则 \(\sin B=\fillinblank{}\).

\begin{answer}
\(\frac{\sqrt{15}}{4}\)
\end{answer}

\begin{solution}
由余弦定理，\(c^2=a^2+b^2-2ab\cos C=1^2+2^2-2\times1\times2\times\frac14=4\)，所以 \(c=2\). 又 \(0<C<\pi\)，故 \(\sin C=\sqrt{1-\cos^2C}=\sqrt{1-\frac1{16}}=\frac{\sqrt{15}}4\). 由正弦定理 \(\frac{b}{\sin B}=\frac{c}{\sin C}\)，且 \(b=c=2\)，得到 \(\sin B=\sin C=\frac{\sqrt{15}}4\).
\end{solution}
\end{problem}

\begin{problem}
设 \(P\) 为直线 \(y=\frac{b}{3a}x\) 与双曲线 \(\frac{x^2}{a^2}-\frac{y^2}{b^2}=1\) （\(a>0\)，\(b>0\)）左支的交点，\(F_1\) 是左焦点，\(PF_1\) 垂直于 \(x\) 轴，则双曲线的离心率 \(e=\fillinblank{}\).

\begin{answer}
\(\frac{3\sqrt{2}}{4}\)
\end{answer}

\begin{solution}
设双曲线的半焦距为 \(c\)，则 \(c^2=a^2+b^2\)，左焦点为 \(F_1=(-c,0)\). 因为 \(PF_1\perp x\) 轴，点 \(P\) 与 \(F_1\) 的横坐标相同，所以 \(x_P=-c\). 由直线方程，\(y_P=-\frac{bc}{3a}\). 将 \(P\) 代入双曲线方程，得 \(\frac{c^2}{a^2}-\frac{c^2}{9a^2}=1\)，即 \(\frac89\frac{c^2}{a^2}=1\). 因此离心率 \(e=\frac ca=\sqrt{\frac98}=\frac{3\sqrt2}{4}\).
\end{solution}
\end{problem}

\begin{problem}
某艺校在一天的 \(6\) 节课中随机安排语文，数学，外语三门文化课和其他三门艺术课各 \(1\) 节，则在课表上的相邻两节文化课之间至少间隔 \(1\) 节艺术课的概率为\fillinblank{}. （用数字作答）

\begin{answer}
\(\frac{1}{5}\)
\end{answer}

\begin{solution}
把语文、数学、外语和三门艺术课视为六门彼此不同的课，课表的总排法数为 \(6!\). 先排三门艺术课，有 \(3!\) 种排法；三门艺术课排好后形成左端、两课之间和右端共 \(4\) 个空隙. 为了使相邻两节文化课之间至少间隔一节艺术课，三门文化课必须分别放入其中三个不同空隙，有 \(\mathrm C_4^3\) 种选法，再排列三门文化课，有 \(3!\) 种排法. 因此有利排法数为 \(3!\mathrm C_4^3 3!\)，所求概率为 \(\frac{3!\mathrm C_4^3 3!}{6!}=\frac15\).
\end{solution}
\end{problem}

\section{解答题}

\begin{problem}
已知 \(\{a_{n}\}\) 为等差数列，且 \(a_{1}+a_{3}=8\)，\(a_{2}+a_{4}=12\).

\begin{enumerate}
\item 求 \(\{a_{n}\}\) 的通项公式；
\item 记 \(\{a_{n}\}\) 的前 \(n\) 项和为 \(\{S_{n}\}\)，若 \(a_{1}\)，\(a_{k}\)，\(a_{k+2}\) 成等比数列，求正整数 \(k\) 的值.
\end{enumerate}
\end{problem}

\begin{solution}
\begin{enumerate}
\item 设等差数列 \(\{a_n\}\) 的公差为 \(d\)，则 \(a_n=a_1+(n-1)d\). 由已知，\(a_1+a_3=2a_1+2d=8\)，所以 \(a_1+d=4\). 又 \(a_2+a_4=2a_1+4d=12\)，于是 \(2(a_1+d)+2d=12\)，解得 \(d=2\)，从而 \(a_1=2\). 因此 \(a_n=2+2(n-1)=2n\).
\item 由前一问，\(a_k=2k\)，\(a_{k+2}=2k+4\). 因为 \(a_1\)，\(a_k\)，\(a_{k+2}\) 成等比数列，所以 \(a_k^2=a_1a_{k+2}\)，即 \((2k)^2=2(2k+4)\). 整理得 \(k^2-k-2=0\)，所以 \((k-2)(k+1)=0\). 由于 \(k\) 为正整数，故 \(k=2\). 此时三项为 \(2\)，\(4\)，\(8\)，确实成等比数列.
\end{enumerate}
\end{solution}

\begin{problem}
已知函数 \(f(x)=ax^3+bx+c\) 在点 \(x=2\) 处取得极值 \(c-16\).

\begin{enumerate}
\item 求 \(a\)，\(b\) 的值；
\item 若 \(f(x)\) 有极大值 28，求 \(f(x)\) 在 \([-3,3]\) 上的最小值.
\end{enumerate}
\end{problem}

\begin{solution}
\begin{enumerate}
\item 由 \(f(2)=c-16\)，得 \(8a+2b+c=c-16\)，即 \(4a+b=-8\). 因为 \(f(x)\) 在 \(x=2\) 处取得极值，且 \(f\) 在该点可导，由可导函数在内点取得极值的必要条件得 \(f'(2)=0\). 又 \(f'(x)=3ax^2+b\)，所以 \(12a+b=0\). 两式相减得 \(8a=8\)，从而 \(a=1\)，\(b=-12\).
\item 此时 \(f(x)=x^3-12x+c\)，且 \(f'(x)=3x^2-12=3(x+2)(x-2)\). 因此 \(f(x)\) 在 \((-\infty,-2)\) 上递增，在 \((-2,2)\) 上递减，在 \((2,+\infty)\) 上递增，故极大值为 \(f(-2)=c+16\). 由极大值为 \(28\)，得 \(c=12\)，所以 \(f(x)=x^3-12x+12\). 在区间 \([-3,3]\) 上只需比较端点和两个临界点的函数值：\(f(-3)=21\)，\(f(-2)=28\)，\(f(2)=-4\)，\(f(3)=3\). 因而 \(f(x)\) 在 \([-3,3]\) 上的最小值为 \(-4\).
\end{enumerate}
\end{solution}

\begin{problem}
甲，乙两人轮流投篮，每人每次投一球. 约定甲先投且先投中者获胜，一直到有人获胜或每人都已投球 \(3\) 次时投篮结束. 设甲每次投篮投中的概率为 \(\frac{1}{3}\)，乙每次投篮投中的概率为 \(\frac{1}{2}\)，且各次投篮互不影响.

\begin{enumerate}
\item 求乙获胜的概率；
\item 求投篮结束时乙只投了 \(2\) 个球的概率.
\end{enumerate}
\end{problem}

\begin{solution}
\begin{enumerate}
\item 记甲投失和投中的概率分别为 \(\frac{2}{3}\)，\(\frac{1}{3}\)，乙投失和投中的概率分别为 \(\frac{1}{2}\)，\(\frac{1}{2}\). 乙在第一次、第二次、第三次投球时获胜的事件分别为：甲先投失后乙投中；前四次投球中甲、乙、甲、乙均投失后乙投中；前五次投球均投失后乙投中. 因而乙获胜的概率为 \(\frac{2}{3}\times\frac{1}{2}+\frac{2}{3}\times\frac{1}{2}\times\frac{2}{3}\times\frac{1}{2}+\left(\frac{2}{3}\times\frac{1}{2}\right)^2\times\frac{2}{3}\times\frac{1}{2}=\frac{1}{3}+\frac{1}{9}+\frac{1}{27}=\frac{13}{27}\).
\item 投篮结束时乙只投了 \(2\) 个球，结束时刻只能是乙第二次投中或甲第三次投中. 前一种情形要求前两轮中甲、乙、甲均投失，乙第二次投中，概率为 \(\frac{2}{3}\times\frac{1}{2}\times\frac{2}{3}\times\frac{1}{2}=\frac{1}{9}\). 后一种情形要求前四次投球均投失，甲第三次投中，概率为 \(\left(\frac{2}{3}\times\frac{1}{2}\right)^2\times\frac{1}{3}=\frac{1}{27}\). 所求概率为 \(\frac{1}{9}+\frac{1}{27}=\frac{4}{27}\).
\end{enumerate}
\end{solution}

\begin{problem}
设函数 \(f(x)=A\sin(\omega x+\varphi)\)（其中 \(A>0\)，\(\omega>0\)，\(-\pi<\varphi\leqslant \pi\)）在 \(x=\frac{\pi}{6}\) 处取得最大值 2，其图象与 \(x\) 轴的相邻两个交点的距离为 \(\frac{\pi}{2}\).

\begin{enumerate}
\item 求 \(f(x)\) 的解析式；
\item 求函数 \(g(x)=\frac{6\cos^{4}x-\sin^{2}x-1}{f\left(x+\frac{\pi}{6}\right)}\) 的值域.
\end{enumerate}
\end{problem}

\begin{solution}
\begin{enumerate}
\item 因为函数的最大值为 \(A\)，所以 \(A=2\). 正弦函数相邻两个零点对应的自变量之差为 \(\frac{\pi}{\omega}\)，由题意得 \(\frac{\pi}{\omega}=\frac{\pi}{2}\)，故 \(\omega=2\). 又函数在 \(x=\frac{\pi}{6}\) 处取得最大值，所以 \(\frac{\pi}{3}+\varphi=\frac{\pi}{2}+2k\pi\)，即 \(\varphi=\frac{\pi}{6}+2k\pi\). 结合 \(-\pi<\varphi\leqslant\pi\)，得 \(\varphi=\frac{\pi}{6}\). 因此 \(f(x)=2\sin\left(2x+\frac{\pi}{6}\right)\).
\item 先化简分子：\[
6\cos^4x-\sin^2x-1=6\cos^4x+\cos^2x-2=(3\cos^2x+2)(2\cos^2x-1)=(3\cos^2x+2)\cos 2x.
\]

又 \[
f\left(x+\frac{\pi}{6}\right)=2\sin\left(2x+\frac{\pi}{2}\right)=2\cos 2x.
\]

当 \(\cos 2x\ne0\) 时，\(g(x)=\frac{3\cos^2x+2}{2}\). 当 \(\cos 2x=0\) 时原式分母为零，故这些点不在定义域内；而 \(\cos 2x=0\) 等价于 \(\cos^2x=\frac{1}{2}\). 因此 \(\cos^2x\) 的取值为 \([0,1]\) 中除去 \(\frac{1}{2}\) 的所有数，从而 \(g(x)\) 的值域为 \(\left[1,\frac{7}{4}\right)\cup\left(\frac{7}{4},\frac{5}{2}\right]\).
\end{enumerate}
\end{solution}

\begin{problem}
如图，在直三棱柱 \(ABC-A_{1}B_{1}C_{1}\) 中，\(AB=4\)，\(AC=BC=3\)，\(D\) 为 \(AB\) 的中点.

\begin{enumerate}
\item 求异面直线 \(CC_{1}\) 和 \(AB\) 的距离；
\item 若 \(AB_{1}\perp A_{1}C\)，求二面角 \(A_{1}-CD-B_{1}\) 的平面角的余弦值.
\end{enumerate}

\bitmapfigure[max width=0.42\linewidth]{img/2012/chongqing_liberal/q20_fig1.png}
\end{problem}

\begin{solution}
\begin{enumerate}
\item 在底面三角形 \(ABC\) 中，\(D\) 为 \(AB\) 的中点，故 \(AD=DB=2\). 由 \(AC=BC\) 及 \(AD=DB\)，\(CD\perp AB\). 在直角三角形 \(ACD\) 中，\(CD=\sqrt{AC^2-AD^2}=\sqrt{9-4}=\sqrt5\). 又 \(CC_1\perp\) 平面 \(ABC\)，所以 \(CD\perp CC_1\)，而 \(CD\perp AB\). 因此 \(CD\) 是异面直线 \(CC_1\) 和 \(AB\) 的公垂线，所求距离为 \(\sqrt5\).
\item 建立空间直角坐标系，取 \(D\) 为原点，令 \(AB\) 在 \(x\) 轴上，\(CD\) 在 \(y\) 轴上，且令棱柱高为 \(h\). 则 \(A=(-2,0,0)\)，\(B=(2,0,0)\)，\(C=(0,\sqrt5,0)\)，\(A_1=(-2,0,h)\)，\(B_1=(2,0,h)\).

由 \(AB_1\perp A_1C\)，得 \(\overrightarrow{AB_1}=(4,0,h)\)，\(\overrightarrow{A_1C}=(2,\sqrt5,-h)\)，从而 \(\overrightarrow{AB_1}\cdot\overrightarrow{A_1C}=8-h^2=0\). 所以 \(h^2=8\). 由于 \(\overrightarrow{DA_1}=(-2,0,h)\) 和 \(\overrightarrow{DB_1}=(2,0,h)\) 都垂直于棱 \(CD\)，故 \(\angle A_1DB_1\) 是二面角 \(A_1-CD-B_1\) 的平面角. 因而其余弦值为 \(\cos\angle A_1DB_1=\frac{\overrightarrow{DA_1}\cdot\overrightarrow{DB_1}}{|DA_1||DB_1|}=\frac{h^2-4}{\sqrt{h^2+4}\sqrt{h^2+4}}=\frac{8-4}{12}=\frac13\).
\end{enumerate}
\end{solution}

\begin{problem}
如图，设椭圆的中心为原点 \(O\)，长轴在 \(x\) 轴上，上顶点为 \(A\)，左，右焦点分别为 \(F_{1}\)，\(F_{2}\)，线段 \(OF_{1}\)，\(OF_{2}\) 的中点分别为 \(B_{1}\)，\(B_{2}\)，且 \(\triangle AB_{1}B_{2}\) 是面积为 \(4\) 的直角三角形.

\begin{enumerate}
\item 求该椭圆的离心率和标准方程；
\item 过 \(B_{1}\) 作直线交椭圆于 \(P\)，\(Q\) 两点，使 \(PB_{2}\perp QB_{2}\)，求 \(\triangle PB_{2}Q\) 的面积.
\end{enumerate}

\bitmapfigure[max width=0.42\linewidth]{img/2012/chongqing_liberal/q21_fig1.png}
\end{problem}

\begin{solution}
\begin{enumerate}
\item 设椭圆的标准方程为 \(\frac{x^2}{a^2}+\frac{y^2}{b^2}=1\)，其中 \(a>b>0\)，并设半焦距为 \(c\)，则 \(c^2=a^2-b^2\). 于是 \(A=(0,b)\)，\(B_1=\left(-\frac{c}{2},0\right)\)，\(B_2=\left(\frac{c}{2},0\right)\). 在点 \(A\) 处，\(\overrightarrow{AB_1}\cdot\overrightarrow{AB_2}=b^2-\frac{c^2}{4}\). 而在点 \(B_1\) 处，\(\overrightarrow{B_1A}\cdot\overrightarrow{B_1B_2}=\frac{c^2}{2}\ne0\)，点 \(B_2\) 处同理不可能为直角，所以直角在点 \(A\) 处. 因此 \(c=2b\). 又 \(4=\frac12|AB_1||AB_2|=\frac12\left(b^2+\frac{c^2}{4}\right)=b^2\)，得 \(b=2\)，\(c=4\)，从而 \(a^2=b^2+c^2=20\). 所以椭圆的离心率为 \(e=\frac ca=\frac{2\sqrt5}{5}\)，标准方程为 \(\frac{x^2}{20}+\frac{y^2}{4}=1\).
\item 由上问，\(B_1=(-2,0)\)，\(B_2=(2,0)\). 过 \(B_1\) 的直线若为竖直直线 \(x=-2\)，其与椭圆交点的纵坐标为 \(\pm\frac{4}{\sqrt5}\)，于是从 \(B_2\) 指向两交点的向量内积为 \(16-\frac{16}{5}\ne0\)，不满足垂直条件. 因此可设该直线为 \(y=m(x+2)\). 令 \(t=x+2\)，则直线上点的坐标为 \((-2+t,mt)\). 代入椭圆方程，交点对应的两个参数 \(t_1\)，\(t_2\) 是方程 \(\frac{(t-2)^2}{20}+\frac{m^2t^2}{4}=1\)，即 \((1+5m^2)t^2-4t-16=0\) 的两个根，故由韦达定理，\(t_1+t_2=\frac{4}{1+5m^2}\)，\(t_1t_2=-\frac{16}{1+5m^2}\). 又 \(\overrightarrow{B_2P}=(-4+t_1,mt_1)\)，\(\overrightarrow{B_2Q}=(-4+t_2,mt_2)\)，由 \(PB_2\perp QB_2\) 得 \(16-4(t_1+t_2)+(1+m^2)t_1t_2=0\). 代入韦达定理的结果，得 \(16-\frac{16}{1+5m^2}-\frac{16(1+m^2)}{1+5m^2}=0\)，所以 \(1+5m^2=2+m^2\)，即 \(m^2=\frac14\). 于是 \(|t_1-t_2|=\frac{\sqrt{16+64(1+5m^2)}}{1+5m^2}=\frac{4\sqrt{10}}{9/4}=\frac{16\sqrt{10}}9\). 最后，\(2S_{\triangle PB_2Q}=\left|(-4+t_1)mt_2-mt_1(-4+t_2)\right|=4|m||t_1-t_2|\)，所以 \(S_{\triangle PB_2Q}=2|m||t_1-t_2|=\frac{16\sqrt{10}}9\).
\end{enumerate}
\end{solution}
